\documentclass{article}
\usepackage{amsmath,amssymb,amsthm}

\begin{document}

\textbf{Maximum-norm a posteriori bound for backward Euler/FEM via Green representation.}

Let u denote the (weak) solution of \(\partial_t u + Mu = F\) on \(\Omega\times(0,T]\) with \(u(0)=u_0\) and \(u|_{\partial\Omega}=0\). Let \(\{u_h^j\}_{j=0}^M\subset V_h^0\) be the backward Euler/FEM approximations, i.e. for \(j=1,\dots,M\)
\[
\delta_t u_h^j\,\cdot_h + c_h(u_h^j,v_h)= (f^j,v_h)_h\quad\forall v_h\in V_h^0,
\]
where \(\delta_t u_h^j=(u_h^j-u_h^{j-1})/\tau_j\), and \(f^j\) is the chosen time-discrete right-hand side on slab endpoints. Define the time extension \(\tilde u_h\) on \((0,T]\) by \(\tilde u_h(t)=\,u_h^j\) for \(t\in I_j\) (or, more generally, by the piecewise affine reconstruction described in the statement). We assume, as in the problem setup for the lemma and the error representation, that \(\tilde u_h\in L^2(0,T;V)\) and that the distributional time derivative \(\partial_t\tilde u_h\in L^2(0,T;V^*)\) is well-defined (with possible jump contributions encoded in the distribution). Set the error
\[
e(t):=u(t)-\tilde u_h(t),\qquad t\in(0,T].
\]
Then \(e\) satisfies homogeneous Dirichlet boundary conditions on \(\partial\Omega\) in the trace sense because both \(u\) and \(\tilde u_h\) take values in \(V=H_0^1(\Omega)\) for each time.

Fix any \(v\in V\). The weak formulations of u and \(\tilde u_h\) in time yield, after subtracting, an identity in \(V^*\)-valued distributions on \((0,T]\). Concretely, since u satisfies
\[
\frac{\mathrm d}{\mathrm dt}(u(t),v)_H + c(u(t),v)=\langle F(t),v\rangle\quad\text{in }\mathcal D'(0,T]\text{ for all }v\in V,
\]
we may write \(\langle Ku(t),v\rangle= \langle F(t),v\rangle\) in distributional sense, with \(Ku=\partial_t u+Mu\). For \(\tilde u_h\), the reconstruction was chosen so that the distributional pairing of \(\partial_t\tilde u_h\) against \(v\) produces the standard temporal defect terms involving \(\delta_t u_h^j\) and the interpolation defects on each slab; moreover, the spatial part uses the elliptic bilinear form with a consistent splitting into the discrete spatial residual \(\psi_h^j\). Recall the discrete elliptic residual element \(\psi_h^j\in V_h^0\) is defined by
\[
(\psi_h^j,v_h)_h:=c_h(u_h^j,v_h)_h-(f^j,v_h)_h\quad\forall v_h\in V_h^0.
\]
This allows one to define, for each slab index \(j\), an elliptic reconstruction \(R^j\in H_0^1(\Omega)\) by
\[
c(R^j,v)= (f^j+\psi_h^j,v)\quad\forall v\in V.
\]
With these definitions, one checks (using only linearity and the definition of \(R^j\)) that the difference \(\tilde u_h(t)-u_h^j\) (for the chosen time reconstruction) accounts for the temporal defect, while \(R^j-u_h^j\) accounts for the spatial/elliptic defect. In particular there exists a distribution \(\rho\in \mathcal D'(0,T;V^*)\) such that
\[
Ke(t)=\rho(t)\quad\text{in }V^*\text{-valued distributions on }(0,T].
\]
Moreover, choosing the time reconstruction exactly as in the statement (piecewise affine using \(u^j\), \(\delta_t u^j\), and midpoint values where applicable), the distribution \(\rho\) can be decomposed into a temporal part and an elliptic part so that for each \(s\in I_j\) and each \(v\in V\), the duality pairing is of the form
\[
\langle \rho(s),v\rangle = \langle \rho_j^{\mathrm{temp}}(s),v\rangle + \langle \rho_j^{\mathrm{ell}},v\rangle,
\]
where \(\rho_j^{\mathrm{ell}}\) depends only on the elliptic reconstruction defect through
\[
\langle \rho_j^{\mathrm{ell}},v\rangle := c(R^j-u_h^j,v),
\]
and \(\rho_j^{\mathrm{temp}}(s)\) is built only from the temporal reconstruction defect and the time-discrete data, namely from \(\delta_t u_h^j\), the chosen slabwise interpolation, and the difference between the exact-in-time right-hand side and the slabwise constant approximation \(f^j\) (as encoded in the residual terms used in the estimator of the problem statement). In particular, there are no additional unknown continuous quantities beyond those already appearing through these reconstruction defects.

Now apply Lemma \(\ref{lem:w-green}\) to the function \(w=e\). The lemma provides, for each fixed \(x\in\Omega\), a representation formula in time for the evaluation \(e(x,t)\) in terms of the initial data and the action of \(Ke\). Since \(e\) has homogeneous boundary conditions (so the lemma applies on \(\Omega\)) and since \(Ke=\rho\) holds in the distributional sense, the lemma yields for all \(t\in(0,T]\) the identity
\[
e(x,t)=\big(\Gamma(t),e(\cdot,0)\big)+\int_0^t \big\langle Ke(s),\Gamma(t-s)\big\rangle\,\mathrm ds
=\big(\Gamma(t),e(\cdot,0)\big)+\int_0^t \big\langle \rho(s),\Gamma(t-s)\big\rangle\,\mathrm ds.
\]
Because \(u_0\) is Hölder continuous and \(u_0|_{\partial\Omega}=0\), the exact solution u admits a continuous representative up to \(\bar\Omega\times[0,T]\), and the Green function estimates together with the standard mapping properties of \(\Gamma(\tau)\in V\) guarantee that the integral representation defines \(e(x,t)\) pointwise for each \(x\) and \(t\). In particular, at the target time \(t=T\) we have
\[
e(x,T)=\big(\Gamma(T),e(\cdot,0)\big)+\int_0^T \big\langle \rho(s),\Gamma(T-s)\big\rangle\,\mathrm ds.
\]
Taking absolute values and then the supremum over \(x\in\Omega\) gives
\[
\|e(T)\|_{\infty,\Omega}\le \big\|\Gamma(T)\big\|_{1,\Omega}\,\|e(\cdot,0)\|_{\infty,\Omega}
+\int_0^T \left|\big\langle \rho(s),\Gamma(T-s)\big\rangle\right|\mathrm ds.
\]
Using the definition of the \(V^*\times V\) duality and the Green estimate assumption \(\|\Gamma(\tau)\|_{1,\Omega}\le \kappa_0 e^{-\gamma\tau}\), we obtain the pointwise bound
\[
\left|\big\langle \rho(s),\Gamma(T-s)\big\rangle\right|\le \|\rho(s)\|_{V^*}\,\|\Gamma(T-s)\|_{1,\Omega}\le \|\rho(s)\|_{V^*}\,\kappa_0 e^{-\gamma(T-s)}.
\]
Therefore
\[
\|u(T)-\tilde u_h(T)\|_{\infty,\Omega}=\|e(T)\|_{\infty,\Omega}
\le \big\|\Gamma(T)\big\|_{1,\Omega}\,\|e(\cdot,0)\|_{\infty,\Omega}
+\kappa_0\int_0^T e^{-\gamma(T-s)}\|\rho(s)\|_{V^*}\,\mathrm ds.
\]
The elliptic/temporal decomposition of \(\rho\) and the restriction \(s\in I_j\) yield a slabwise decomposition of the integral. Indeed, writing \(\int_0^T = \sum_{j=1}^M\int_{I_j}\) and defining the slab weights
\[
\omega_j:=\kappa_0\int_{I_j} e^{-\gamma(T-s)}\,\mathrm ds,
\]
we infer
\[
\|e(T)\|_{\infty,\Omega}\le \big\|\Gamma(T)\big\|_{1,\Omega}\,\|e(\cdot,0)\|_{\infty,\Omega}
+\sum_{j=1}^M \omega_j\sup_{s\in I_j}\|\rho(s)\|_{V^*}.
\]
To make this explicit in terms of computable quantities, we evaluate the duality pairings using the two parts of \(\rho\). For each slab \(I_j\), the elliptic part contributes
\[
\int_{I_j}\left|\big\langle \rho_j^{\mathrm{ell}},\Gamma(T-s)\big\rangle\right|\mathrm ds
=\int_{I_j}\left|c(R^j-u_h^j,\Gamma(T-s))\right|\mathrm ds.
\]
By continuity of c with respect to the \(H^1\)-norm, and since \(\Gamma(T-s)\in V\) with \(\|\Gamma(T-s)\|_{1,\Omega}\) controlled by Assumption \(\ref{ass:green}\), we obtain
\[
\left|c(R^j-u_h^j,\Gamma(T-s))\right|\le C\|R^j-u_h^j\|_{\infty,\Omega}\,\|\Gamma(T-s)\|_{1,\Omega}.
\]
Using Assumption \(\ref{ass:ell-est}\) for the elliptic problem underlying the reconstruction, applied to the continuous solution \(R^j\) and the discrete solution \(u_h^j\) with right-hand side \(f^j+\psi_h^j\), we get
\[
\|R^j-u_h^j\|_{\infty,\Omega}\le \eta_\mathrm{ell}^j,
\]
where \(\eta_\mathrm{ell}^j:=\eta_\mathrm{ell}(u_h^j,f^j+\psi_h^j)\) is exactly the elliptic estimator defined in the statement.
Consequently the elliptic contribution over \(I_j\) is bounded by a computable multiple of \(\eta_\mathrm{ell}^j\), with the multiplicative factor depending only on the Green bounds (hence only on \(\kappa_0,\gamma\) and \(T\)).

For the temporal part, \(\rho_j^{\mathrm{temp}}(s)\) is built from the defect between the continuous time evolution of \(\tilde u_h\) and the backward Euler/FEM law; with the chosen slabwise reconstruction this defect is a linear functional of quantities of the form \(\delta_t u_h^j\) and the interpolation/reconstruction defects (for example the difference between the slab midpoint reconstruction and the actual piecewise affine representation, exactly as in the problem’s definition of the estimator). Hence, for each \(s\in I_j\) the duality pairing \(\langle \rho_j^{\mathrm{temp}}(s),\Gamma(T-s)\rangle\) is bounded by a constant times \(\|\Gamma(T-s)\|_{1,\Omega}\) times a computable norm of these temporal defect functionals. By Assumption \(\ref{ass:green}\), \(\|\Gamma(T-s)\|_{1,\Omega}\) is controlled by \(\kappa_0 e^{-\gamma(T-s)}\) (and \(\|\partial_t\Gamma\|_{1,\Omega}\) is used if the derivation of the temporal defect representation includes integration by parts in time; in that case the same assumption provides the bound through the Green estimate for \(\partial_t\Gamma\)). Therefore each temporal contribution on \(I_j\) is bounded by a computable weight (depending only on \(\kappa_0,\gamma,\kappa_1,\kappa_1'\) as needed) times the temporal estimator term involving \(\delta_t u_h^j\) and the reconstruction/interpolation defects.

Finally, consider the initial term \((\Gamma(T),e(\cdot,0))\). If the extension is chosen so that \(\tilde u_h(0)=u_0\) in the reconstruction (as is compatible with the error analysis framework of the lemma), then \(e(\cdot,0)=0\) and this contribution vanishes. Otherwise, the initial term is bounded analogously by
\[
\big|(\Gamma(T),e(\cdot,0))\big|\le \|\Gamma(T)\|_{1,\Omega}\,\|e(\cdot,0)\|_{\infty,\Omega},
\]
which is incorporated as an additional computable term in the global estimator, based on the given initial data mismatch.

Collecting the initial term and the slabwise temporal and elliptic contributions, and using only that \(u_h^\sim(T)=u_h(T)\) by construction of the time extension at the final time node, we arrive at an inequality of the form
\[
\|u(T)-u_h(T)\|_{\infty,\Omega}\le \eta_{\max}\big(u_h,\{\psi_h^j\}_{j=1}^M,\{\delta_t u_h^j\}_{j=1}^M,\{\eta_\mathrm{ell}^j\}_{j=1}^M\big),
\]
where \(\eta_{\max}\) is defined as the sum over slabs of the temporal defect indicators multiplied by the Green weights and the elliptic indicators \(\eta_\mathrm{ell}^j\) multiplied by the corresponding Green weights, plus the initial-data term (if present). By construction from the representation identity and the preceding bounds, \(\eta_{\max}\) majorizes exactly the right-hand side produced by Lemma \(\ref{lem:w-green}\), so it is a maximum-norm a posteriori upper bound for the error at time \(T\) in the backward Euler/FEM discretization.

This completes the derivation of a maximum-norm a posteriori estimator for \(\|u(T)-u_h(T)\|_{\infty,\Omega}\).

\end{document}